% !TEX root = ../notes_template.tex
\chapter{Shoulder Complex}\label{chp:shoulder_complex}

\section*{Objectives}
\begin{itemize}
  \item Perform and interpret manual muscle testing (MMT) of glenohumeral and scapulothoracic musculature.
  \item Identify each tested muscle’s innervation and most common spinal root(s).
  \item Conduct focused muscle length testing (MLT) of upper quarter muscles relevant to shoulder mechanics.
  \item Palpate and image the axillary artery using ultrasound.
  \item Acquire and interpret static and dynamic rehabilitative ultrasound images of the supraspinatus tendon and AC region.
\end{itemize}

\section*{Functional Context and Movement Relevance}

The shoulder is a highly mobile region that depends on coordinated scapulohumeral rhythm for functional movement. Proper glenohumeral motion is dependent on the ability of the scapula to rotate upward to maintain congruence between the humeral head and the glenoid fossa. Dysfunction of the scapulothoracic joint often manifests itself as impaired shoulder strength, aberrant movement patterns, or compensatory postures.

This week builds directly on the functional movement lens of Week 1 by focusing on the Selective Functional Movement Assessment (SFMA) upper extremity top-tier patterns:

\begin{itemize}
  \item \textbf{UE1:} Shoulder internal rotation with extension – reach behind back.
  \item \textbf{UE2:} Shoulder external rotation with abduction – reach behind head.
\end{itemize}

These movements help reveal asymmetries, motion restrictions, and potential motor control deficits in the shoulder complex. In this lab, you'll build from these patterns to investigate glenohumeral and scapular muscle function more precisely through MMT, MLT, and ultrasound-based evaluation.

\section*{Shoulder Complex Muscles by Attachment}

\noindent
Understanding shoulder complex muscles by their anatomical attachments—whether they span from the thorax, scapula, or clavicle—helps clarify how they contribute to scapular positioning and glenohumeral motion. This organization is not only anatomically meaningful, but functionally essential when performing manual muscle testing (MMT). 

Knowing which region anchors the muscle informs both \textbf{where stabilization should occur} and \textbf{how resistance should be applied} to bias the target muscle. For example, testing a thoracoscapular muscle like the rhomboids requires stabilizing the thorax and observing scapular movement, whereas testing a scapulohumeral muscle like the deltoid focuses on glenohumeral motion with attention to scapular positioning. This attachment-based framework supports more accurate test setup, interpretation, and clinical reasoning.

\vspace{.5cm}

\begin{tabular}{|p{4.5cm}|p{4.5cm}|p{5cm}|}
\hline
\textbf{Thoracoscapular} & \textbf{Thoracohumeral} & \textbf{Scapulohumeral} \\
\hline
Trapezius &
Latissimus Dorsi &
Deltoid  \\
Serratus Anterior &
Pectoralis Major &
Supraspinatus \\
Rhomboid Major &
 &
Infraspinatus \\
Rhomboid Minor &
 &
Subscapularis \\
Levator Scapulae &
 &
Teres Minor \\
Pectoralis Minor &
 &
Teres Major \\
Subclavius &
 &
Biceps Brachii  \\
 &
 &
Triceps Brachii (long head) \\
\hline
\end{tabular}

\section*{Scapular Movements: Muscles and Innervation}

\vspace{.5cm}

\noindent
The table below organizes scapular osteokinematic movement patterns by their primary muscular contributors (synergists for the motion), along with each muscle’s peripheral nerve and most common nerve root(s). Muscles are grouped by movement and listed top-down by ascending nerve root to aid in clinical reasoning and motor screening.

\begin{tabular}{|p{4cm}|p{4.5cm}|p{4cm}|p{2.5cm}|}
\hline
\textbf{Scapular Motion} & \textbf{Muscle} & \textbf{Peripheral Nerve} & \textbf{Nerve Root(s)} \\
\hline
\multirow{3}{*}{Elevation} 
  & Levator Scapulae & Dorsal Scapular & C3–C5 \\
  & Upper Trapezius & Spinal Accessory & CN XI, C1–C5 \\
  & Rhomboid Major & Dorsal Scapular & C4–C5 \\
\hline
\multirow{3}{*}{Depression}
  & Pectoralis Minor & Medial Pectoral & C8–T1 \\
  & Latissimus Dorsi & Thoracodorsal & C6–C8 \\
  & Lower Trapezius & Spinal Accessory & CN XI, C1–C5 \\
\hline
\multirow{3}{*}{Upward Rotation}
  & Serratus Anterior & Long Thoracic & C5–C7 \\
  & Upper Trapezius & Spinal Accessory & CN XI, C1–C5 \\
  & Lower Trapezius & Spinal Accessory & CN XI, C1–C5 \\
\hline
\multirow{4}{*}{Downward Rotation}
  & Levator Scapulae & Dorsal Scapular & C3–C5 \\
  & Rhomboid Minor & Dorsal Scapular & C4–C5 \\
  & Rhomboid Major & Dorsal Scapular & C4–C5 \\
  & Pectoralis Minor & Medial Pectoral & C8–T1 \\
\hline
\multirow{2}{*}{Retraction}
  & Rhomboid Major & Dorsal Scapular & C4–C5 \\
  & Middle Trapezius & Spinal Accessory & CN XI, C1–C5 \\
\hline
\multirow{2}{*}{Protraction}
  & Serratus Anterior & Long Thoracic & C5–C7 \\
  & Pectoralis Minor & Medial Pectoral & C8–T1 \\
\hline
\end{tabular}

\begin{tcolorbox}[colback=gray!5, colframe=black, title={Reminder: Movement, Not Muscle Isolation}]
Manual muscle testing of the shoulder complex emphasizes movements, not isolated muscles. While positioning and resistance can bias certain muscles, remember that each test reflects the performance of a movement system involving synergists, stabilizers, and coordinated motor control. Use the tables above to inform your testing, but interpret findings within the broader functional context of scapulohumeral rhythm and regional interdependence.
\end{tcolorbox}

\begin{labactivitybox}
\textbf{Activity: Scapular MMT}
\begin{itemize}
  \item Test scapular elevation, depression, protraction, retraction, and upward/downward rotation.
  \item Document observed compensation and stabilization strategies.
  \item Identify peripheral nerve and root level for each primary muscle tested.
\end{itemize}
\end{labactivitybox}

\begin{labactivitybox}
\textbf{Activity: Muscle Length Testing of Scapulothoracic Influencers}
\begin{itemize}
  \item Perform MLT for: Pectoralis Major (sternal and clavicular heads), and Pectoralis Minor.
  \item Document R1 and R2 (points of initial and firm resistance or observed compensation).
  \item Classify each result as Limited, Normal, or Excessive based on symmetry and expected range.
\end{itemize}
\end{labactivitybox}

\section*{Glenohumeral Movements: Muscles and Innervation}

\vspace{.5cm}

\noindent
The table below organizes glenohumeral (GH) movement patterns by their primary muscular contributors, along with each muscle’s peripheral nerve and most common nerve root(s). Muscles are grouped by movement and listed top-down by ascending nerve root to support neuroanatomical reasoning.

\begin{table}[H]
\centering
\caption{Glenohumeral Joint Movements, Prime Movers, and Innervation}
\begin{tabular}{|p{4cm}|p{4.5cm}|p{4cm}|p{2.5cm}|}
\hline
\textbf{GH Movement} & \textbf{Muscle} & \textbf{Peripheral Nerve} & \textbf{Nerve Root(s)} \\
\hline
\multirow{3}{*}{Flexion}
  & Biceps Brachii & Musculocutaneous & C5–C6 \\
  & Anterior Deltoid & Axillary & C5–C6 \\
  & Coracobrachialis & Musculocutaneous & C5–C7 \\
\hline
\multirow{4}{*}{Extension}
  & Teres Major & Subscapular (low) & C5–C6 \\
  & Posterior Deltoid & Axillary & C5–C6 \\
  & Latissimus Dorsi & Thoracodorsal & C6–C8 \\
  & Triceps Brachii (lh) & Radial & C6–C8 \\
\hline
\multirow{2}{*}{Abduction}
  & Supraspinatus & Suprascapular & C5–C6 \\
  & Middle Deltoid & Axillary & C5–C6 \\
\hline
\multirow{3}{*}{Adduction}
  & Teres Major & Subscapular (low) & C5–C6 \\
  & Latissimus Dorsi & Thoracodorsal & C6–C8 \\
  & Pectoralis Major & Pectoral (med \& lat) & C5–T1 \\
\hline
\multirow{3}{*}{External Rotation}
  & Infraspinatus & Suprascapular & C5–C6 \\
  & Teres Minor & Axillary & C5–C6 \\
  & Supraspinatus\textsuperscript{1} & Suprascapular & C5–C6 \\
\hline
\multirow{4}{*}{Internal Rotation}
  & Subscapularis & Subscapular (up \& low) & C5–C6 \\
  & Teres Major & Subscapular (low) & C5–C6 \\
  & Latissimus Dorsi & Thoracodorsal & C6–C8 \\
  & Pectoralis Major & Pectoral (med \& lat) & C5–T1 \\
\hline
\end{tabular}
\end{table}

\footnotetext[1]{The supraspinatus is not a primary external rotator, but may contribute minimally to ER torque in some positions. Its primary functions are initiating abduction and compressing the humeral head. Clinically, it is often recruited during ER-focused rehab tasks due to its role in joint stability, not due to direct ER torque production.}

\begin{tcolorbox}[colback=gray!5, colframe=black, title={Final Reminder: Movement, Not Muscle Isolation}]
Manual muscle testing of the shoulder complex emphasizes movements, not isolated muscles. While positioning and resistance can bias certain muscles, remember that each test reflects the performance of a movement system involving synergists, stabilizers, and coordinated motor control. Use the tables above to inform your testing, but interpret findings within the broader functional context of scapulohumeral rhythm and regional interdependence.
\end{tcolorbox}


\begin{labactivitybox}
\textbf{Activity: Glenohumeral MMT}
\begin{itemize}
  \item Perform gravity-resisted and gravity-eliminated testing for the GH motions.
  \item Apply break test and record HHD output for one motion.
  \item Identify innervation and nerve root(s) for each muscle tested.
\end{itemize}
\end{labactivitybox}

\begin{labactivitybox}
\textbf{Activity: Muscle Length Testing of GH Influencers}
\begin{itemize}
  \item Perform muscle length tests for: Pectoralis Major (sternal/clavicular), Pectoralis Minor, Latissimus Dorsi, and Teres Major.
  \item Record R1 and R2 compensations and classify each as Limited, Normal, or Excessive.
\end{itemize}
\end{labactivitybox}


\section*{Rehabilitative Ultrasound Imaging (RUSI) of the Glenohumeral Complex}

Ultrasound imaging of the shoulder offers a unique opportunity to visualize superficial structures relevant to both movement and clinical dysfunction. In this lab, RUSI is used not to diagnose pathology, but to:

\begin{itemize}
  \item Reinforce layered anatomical relationships through real-time visualization.
  \item Practice probe handling, alignment, and anisotropy correction in a high-utility region.
  \item Identify landmarks and dynamic tissue interactions during functional shoulder motion.
\end{itemize}

The glenohumeral joint is a prime region for developing RUSI skills due to the accessibility of key tendons, the clinical relevance of shoulder dysfunction, and the well-established sonographic windows for the biceps tendon and supraspinatus outlet.

The following scans focus on clinically useful structures:
\begin{itemize}
  \item The long head of the biceps tendon — a foundational scan for shoulder RUSI orientation.
  \item The supraspinatus tendon in SAX and LAX — for assessment of tendon integrity and depth relationships.
  \item A dynamic view of the subacromial space during abduction — to explore potential impingement mechanisms.
\end{itemize}

\footnotetext{Adapted from: Collebrusco et al. (2017), \textit{The Use of Ultrasound Images in Manual Therapy}; and Bailey et al. (2015), \textit{Current Rehabilitation Applications for Shoulder Ultrasound Imaging}.}

\begin{labactivitybox}
\textbf{Activity: Biceps Tendon (Long Head) – Revisited for Shoulder Orientation}

This scan revisits the long head of the biceps tendon as a foundational reference for shoulder RUSI. Its central location within the bicipital groove provides a key landmark for orienting future scans of the rotator cuff and shoulder joint.

\textit{Steps:}
\begin{enumerate}
  \item Position the patient seated, arm relaxed and slightly externally rotated.
  \item Palpate and locate the bicipital groove just inferior to the anterior acromion.
  \item Begin with the \textbf{short axis (SAX)} view: the tendon should appear as a well-defined hyperechoic oval bordered by the greater and lesser tubercles.
  \item Rotate the probe 90° to obtain the \textbf{long axis (LAX)} view: observe the continuity of the tendon along its proximal-to-distal path.
  \item Optimize \textbf{anisotropy correction} through heel-toe adjustments, and refine image clarity with gain and depth settings.
\end{enumerate}

\textit{Record:}
\begin{itemize}
  \item What bony landmarks helped confirm SAX and LAX views?
  \item Were tendon borders clearly visualized? If not, how did you adjust probe orientation?
  \item How does this scan reinforce your anatomical model of the anterior shoulder?
  \item How might you use this position to pivot toward supraspinatus or subscapularis imaging?
\end{itemize}
\end{labactivitybox}

\begin{labactivitybox}
\textbf{Activity: Supraspinatus Tendon Imaging (Static and Dynamic)}

This scan introduces the supraspinatus tendon and the subacromial space in both SAX and LAX views, highlighting its relationship to the humeral head, greater tuberosity, acromion, and subdeltoid bursa. It includes a dynamic screen for anterior impingement.

\textit{Steps (Static Imaging):}
\begin{enumerate}
  \item Position the patient seated with the hand behind the back or on the thigh, shoulder slightly extended.
  \item \textbf{SAX View:} Place the probe transversely just anterolateral to the acromion to visualize the convex humeral head (red outline), overlying supraspinatus tendon (“tire” appearance), and subdeltoid bursa (black hypoechoic interface).
  \item \textbf{LAX View:} Rotate the probe obliquely to follow the tendon along the greater tuberosity to the musculotendinous junction. Identify a continuous, beak-shaped tendon with uniform echotexture. 
  \item Confirm tendon integrity and bursal separation from the deltoid muscle. Adjust for anisotropy.
\end{enumerate}

\textit{Steps (Dynamic Imaging for Impingement):}
\begin{enumerate}
  \item Maintain LAX orientation and passively abduct the arm from ~30° to 90°.
  \item Observe whether the supraspinatus tendon, subdeltoid bursa, or overlying deltoid compresses or narrows the acromiohumeral space.
  \item If available, use cine capture to analyze any “shearing” of the bursa during elevation.
\end{enumerate}

\textit{Record:}
\begin{itemize}
  \item Describe tendon appearance and continuity in both SAX and LAX views.
  \item Was subacromial spacing preserved or narrowed during abduction?
  \item Identify any signs of bursal impingement or deltoid encroachment.
\end{itemize}
\end{labactivitybox}

\begin{labactivitybox}
\textbf{Activity: Dynamic AC Joint Imaging – Screening for Subacromial Impingement}

This scan assesses the supraspinatus outlet and subacromial space during passive abduction to identify or rule out dynamic impingement.

\textit{Steps:}
\begin{enumerate}
  \item Place the probe in an \textbf{oblique long axis (LAX)} view over the supraspinatus tendon just inferior to the acromion.
  \item Identify the \textbf{acromion (Acr)}, \textbf{humeral head (Hum)}, and the overlying \textbf{subacromial-subdeltoid bursa}.
  \item Actively abduct the shoulder through midrange (approx. 30°–90°), keeping the probe stable.
  \item Observe for narrowing of the subacromial space, displacement or “shearing” of the tendon, or bursal effacement.
  \item Use cine loop review to assess motion patterns and tendon glide if available.
\end{enumerate}

\textit{Record:}
\begin{itemize}
  \item Was the anechoic bursal space maintained throughout abduction?
  \item Did the supraspinatus tendon deform, compress, or lose continuity under the acromion?
  \item How do these dynamic findings relate to your static supraspinatus imaging?
  \item Suggested labeling for captured loops: \texttt{Rt Ant Impng Pre} / \texttt{Post Dynamic}
\end{itemize}
\end{labactivitybox}


\section*{Pulse Palpation}

The axillary artery is the principal vascular supply to the upper extremity. As it courses beneath the clavicle and through the thoracic outlet, it is subject to compression by surrounding structures such as the scalene muscles, first rib, and pectoralis minor. This vulnerability contributes to the vascular subtype of thoracic outlet syndrome (TOS), in which mechanical compression compromises distal perfusion and contributes to symptoms such as numbness, coolness, or fatigue during upper extremity activity.

In clinical practice, assessing the axillary pulse provides insight into proximal vascular integrity. Combining palpation with RUSI enhances accuracy and reinforces anatomical understanding.

\begin{labactivitybox}
\textbf{Activity: Pulse Palpation + Imaging}
\begin{itemize}
  \item Locate and palpate the axillary pulse in a relaxed, abducted position.
  \item Image the vessel in short axis using ultrasound; assess for:
    \begin{itemize}
      \item Compressibility (vein collapses, artery remains open).
      \item Pulsatility (real-time rhythmic motion).
      \item Spatial relationships to surrounding muscles (pectoralis major/minor, deltoid).
    \end{itemize}
\end{itemize}
\end{labactivitybox}


\section*{Clinical Integration}

\begin{itemize}
  \item How did observing scapular and glenohumeral motion during movement screening influence your MMT and MLT strategy?
  \item When testing muscles around the shoulder complex, where did you find it difficult to isolate a single muscle? How did your knowledge of innervation help you interpret findings?
  \item How does the concept of regional interdependence apply to abnormal scapulohumeral rhythm or pain during shoulder movement?
  \item What role did ultrasound imaging play in reinforcing your understanding of tendon orientation and joint clearance?
  \item How did your pulse palpation findings support your understanding of the vascular anatomy of the shoulder region?
\end{itemize}

\section*{Next Steps}

\begin{itemize}
  \item Review muscle actions, innervations, and typical compensatory strategies for scapular and glenohumeral motions.
  \item Revisit the ultrasound scanning protocol for the supraspinatus and subacromial space, with attention to dynamic imaging cues.
  \item Practice shoulder MMTs with a peer, focusing on stabilization and resisting substitutions.
  \item Prepare for Week 4 by previewing distal upper extremity movements, including forearm pronation/supination and grip.
\end{itemize}

\section*{Checklist Summary}

\begin{itemize}
  \item [$\square$] Performed gravity-resisted and gravity-eliminated MMT for glenohumeral and scapular muscles.
  \item [$\square$] Identified the innervation and nerve root(s) of all muscles tested during MMT.
  \item [$\square$] Used HHD to record force output for one GH break test.
  \item [$\square$] Conducted and interpreted MLT for pectoralis major, pectoralis minor, latissimus dorsi, and teres major, including R1, R2, and compensations.
  \item [$\square$] Acquired static ultrasound images of the supraspinatus tendon in SAX and LAX.
  \item [$\square$] Performed dynamic ultrasound to assess AC joint space during shoulder abduction.
  \item [$\square$] Palpated and imaged the axillary pulse, confirming vascular characteristics via ultrasound.
\end{itemize}
